\documentclass[10pt,a4paper]{amsart}
\usepackage[margin=19mm]{geometry}
\usepackage{amsmath,amssymb,mathtools}
\usepackage[scr=rsfs,cal=euler]{mathalfa}
\usepackage{xfrac}
\usepackage[unicode,hidelinks,bookmarks=false]{hyperref}
\newcommand{\ie}{i.e.\ }
\newcommand{\cf}{cf.\ }
\newcommand{\resp}{respectively\ }
\newcommand{\Subsection}{\subsection}
\providecommand{\nobreakdash}{\nobreak\hbox{-}}
\setlength{\emergencystretch}{2em}
\multlinegap0pt
\usepackage{amssymb}
\usepackage{mathtools}
\usepackage{xcolor}
\newtheorem{theorem}{Theorem}
\title{A parity Erd\H{o}s--Hajnal theorem for \(t\)-intersecting curves}
\author{Andrew Suk, Su Zhou}
\date{Source: arXiv:2606.11649v1 (CC BY 4.0)}
\begin{document}
\maketitle

\noindent\emph{Source and licence.} A parity Erd\H{o}s--Hajnal theorem for \(t\)-intersecting curves. Version arXiv:2606.11649v1 (CC BY 4.0), distributed by arXiv under the Creative Commons Attribution 4.0 International licence, http://creativecommons.org/licenses/by/4.0/, as recorded in the arXiv licence metadata for this exact version and shown on the abstract page https://arxiv.org/abs/2606.11649v1. Full licence notice: \texttt{licenses/arxiv-2606.11649-CC-BY-4.0.txt}.\par\smallskip

\noindent\textit{Editorial scope.} Counted unit: the article's theorem thm:red-blue-region (source0.tex lines 654--676). The article's complete original proof of it is delivered with it (source0.tex lines 679--802, one \texttt{proof} environment of 124 lines). No mathematics is written, repaired or re-derived: the statement and the proof are the article's own lines, in the article's own notation, and no retained line is substituted or re-flowed.  Retained uncounted as the source-local context the proof needs: lines 319--329.  Nothing was omitted from any retained span, so the package carries no omission record.  No retained line contains a citation, so the wrapper carries no bibliography and no bibliography entry.  No retained line contains a figure environment, an image-inclusion command or drawing code, so no image is delivered and the package declares no asset dependency.  Editorial formatting only, in addition: an AMS \texttt{amsart} preamble that restates the article's own declarations and macros used by the retained lines, namely the environments \texttt{theorem} on separate counters, together with the standard packages those lines need.  The retained environments print in this document as Theorem 1 in order of appearance, whereas the article prints the same items with the numbering of its own full document.  The complete native source of the article as distributed in the arXiv e-print for this exact version is bundled unchanged as \texttt{sources/202564/source0.tex}.
\begin{theorem}\label{thm:mixed}
There exists an absolute constant \(r_0>0\) such that the following holds. Let \(G\) be a triangulated planar graph on \(N\) vertices embedded in the plane, and let \(Q\subseteq V(G)\). Then for every \(r\ge r_0\) and every \(\rho\ge 1\), there exists a division \(\mathcal R\) of \(G\) such that
\[
|\mathcal R|=O\!\left(\frac{N}{r}+\frac{|Q|}{\rho}\right),
\]
and every region \(R\in\mathcal R\) satisfies
\[
|V(R)|=O(r),\qquad |b(R)|=O(\sqrt r),\qquad |(V(R)\setminus b(R))\cap Q|=O(\rho),
\]
and has at most \(24\) holes.
\end{theorem}

\begin{theorem}\label{thm:red-blue-region}
For every \(t\ge 1\), the following holds. Let \(\mathcal B\) be a set of
\(m\) blue curves, and let \(\mathcal G\) be a set of \(n\) green curves, all
contained in a domain \(\Delta\), such that \(\mathcal B\cup\mathcal G\) is
\(t\)-intersecting and in general position.

Then, for every \(\delta\in(0,1)\) and every nonempty point set \(P\)
obtained by selecting at most one endpoint from each curve in \(\mathcal B\),
there is a domain $\Delta'$ with \(\operatorname{cl}(\Delta')\subset\Delta\) such that the following hold.
\begin{enumerate}
    \item $|P\cap\Delta'|\ge \Omega\left({\delta\over t}|P|\right).$

    \item At most \(O(\lceil\sqrt\delta(m+n)\rceil)\) curves in \(\mathcal G\) intersect
    \(\operatorname{cl}(\Delta')\).

    \item The total number of crossing points between the curves in
    \(\mathcal B\cup\mathcal G\) and \(\partial\Delta'\) is
    \(O(\lceil\sqrt\delta(m+n)\rceil)\).

    \item \(\Delta\setminus\operatorname{cl}(\Delta')\) has at most \(24\)
    path-connected components.
\end{enumerate}
\end{theorem}

\begin{proof}
Let $\mathcal B=\{\beta_1,\ldots,\beta_m\}, \mathcal G=\{\gamma_1,\ldots,\gamma_n\}$. We may assume that \(\delta(m+n)^2\ge r_0\), where \(r_0\) is the absolute
constant from Theorem~\ref{thm:mixed}. Otherwise the statement follows by taking \(\Delta'\) to be a sufficiently
small disk around one point of \(P\), chosen in general position.

We construct a plane graph \(G\) as follows. Place a vertex at each
intersection point of two curves in \(\mathcal B\cup\mathcal G\), and also at
the two endpoints of every curve. For each green curve \(\gamma_i\), choose
one additional point \(q_i\in\gamma_i\) that is neither an endpoint nor an
intersection point, and add \(q_i\) as a vertex. Let
\[
Q=\{q_1,\ldots,q_n\}.
\]
The edges of \(G\) are the arcs between consecutive vertices along the curves
in \(\mathcal B\cup\mathcal G\). Thus \(G\) is a plane multigraph contained
in \(\Delta\).   

Since \(\mathcal B\cup\mathcal G\) is \(t\)-intersecting, the number of
intersection points determined by its curves is at most \(t(m+n)^2\). We
refine the plane graph \(G\) inside \(\Delta\) as follows.  Faces and their
sizes are understood in the usual plane embedding of \(G\), with the size of a
face equal to the length of its boundary walk. We eliminate
faces of size \(2\) by, inside $\Delta$, placing one new vertex inside each such face and
joining it to the two boundary vertices. This adds \(O(t(m+n)^2)\) vertices.
We then triangulate all faces of size at least \(4\), again inside
\(\Delta\), and let \(G'\) be the resulting triangulated plane graph. Hence
\[
|V(G')|=O(t(m+n)^2).
\]



Apply Theorem~\ref{thm:mixed} to \(G'\), with the set \(Q\), and with
parameters
\[
r=\delta(m+n)^2,\qquad \rho=\sqrt\delta(m+n).
\]
As a result, we obtain a division \(\mathcal R\) of \(G'\) such that every region
\(R\in\mathcal R\) satisfies
\[
|V(R)|=O(\delta(m+n)^2),\qquad
|b(R)|=O(\sqrt\delta(m+n)),\qquad
|(V(R)\setminus b(R))\cap Q|=O(\sqrt\delta(m+n)),
\]
and has at most \(24\) holes. Moreover,
\[
|\mathcal R|
=
O\left({|V(G')|\over r}+{|Q|\over \rho}\right)
=
O\left({t\over\delta}\right).
\]

For each \(R\in\mathcal R\), let \(\Delta_R\subset \Delta\) be obtained from
the embedded edges of \(R\), together with \(f\cap\Delta\) for every face
\(f\) of \(R\) that is also a face of \(G'\). Since \(R\) is connected,
\(\Delta_R\) is path-connected. Since the regions cover all edges of \(G'\),
the sets \(\Delta_R\) with \(R\in\mathcal R\), cover all vertices of \(G'\), and
hence cover \(P\).  

We next bound the number of path-connected components of
\(\Delta\setminus\Delta_R\). Since \(\Delta_R\) contains the embedded graph
\(R\), each such component is contained in a single face of \(R\). The faces
of \(R\) that are also faces of \(G'\) contribute nothing, by the definition of
\(\Delta_R\). Hence every component of \(\Delta\setminus\Delta_R\) is contained
in \(\Delta\cap h\) for some hole \(h\) of \(R\). Since \(\Delta\) is a domain
and \(\partial h\subset\Delta\), each set \(\Delta\cap h\) is path-connected.
Thus each hole of \(R\) contributes at most one component of
\(\Delta\setminus\Delta_R\). As \(R\) has at most \(24\) holes,
\(\Delta\setminus\Delta_R\) has at most \(24\) path-connected components.




By the pigeonhole principle, there is a region \(R\in\mathcal R\) such that
\[
|P\cap\Delta_R|
\ge { |P| \over |\mathcal R|}
=
\Omega\left({\delta\over t}|P|\right).
\]
Fix this region \(R\).

We next bound the number of green curves that meet \(\Delta_R\). A curve in
\(\mathcal B\cup\mathcal G\) can enter or leave \(\Delta_R\) only through a
boundary vertex of \(R\). By general position, each such boundary vertex lies
on at most two curves. Hence the total number of entry and exit points of all
curves through \(\partial\Delta_R\) is
\[
O(|b(R)|)=O(\sqrt\delta(m+n)).
\]
We claim that
\[
\bigl|\{\gamma_i\in\mathcal G:\gamma_i\cap\Delta_R\ne\emptyset\}\bigr|
\le |(V(R)\setminus b(R))\cap Q|+3|b(R)|.
\]
Indeed, if \(q_i\in V(R)\setminus b(R)\), then \(\gamma_i\) is counted by the
first term. If \(q_i\in b(R)\), then it is counted by the boundary term.
Otherwise \(q_i\notin V(R)\). If \(\gamma_i\) meets \(\Delta_R\), then,
traversing \(\gamma_i\) from \(q_i\) to its first point of entry into
\(\Delta_R\), we meet \(\partial\Delta_R\). Thus \(\gamma_i\) enters
\(\Delta_R\) through a boundary vertex of \(R\). By general position, each
boundary vertex is incident to at most two green curves. This proves the
claim.  By the bounds on \(R\), at most $O(\sqrt\delta(m+n))$ green curves meet \(\Delta_R\).

We may choose a sufficiently small open neighborhood \(\Delta'\) of
\(\Delta_R\), with path-connected closure and
\(\operatorname{cl}(\Delta')\subset\Delta\), chosen so close to \(\Delta_R\)
that \(\Delta\setminus\operatorname{cl}(\Delta')\) has at most \(24\)
path-connected components.  We also choose it small enough so that every green curve
disjoint from \(\Delta_R\) is also disjoint from \(\operatorname{cl}(\Delta')\).
Finally, since $\mathcal B\cup \mathcal G$ is in general position, we can choose \(\partial\Delta'\) so that the number of crossing points
between \(\partial\Delta'\) and the curves in $\mathcal B\cup \mathcal G$ is at most four times the
number of boundary vertices of \(R\). Therefore, the number of crossing points between the curves in $\mathcal B\cup\mathcal G$ and $\partial\Delta'$ is at most $O(\sqrt\delta(m+n))$.  Moreover, 

\[
|P\cap\Delta'|
\ge |P\cap\Delta_R|
\ge \Omega\left({\delta\over t}|P|\right).
\]
Also, by the choice of \(\Delta'\), at most \(O(\sqrt\delta(m+n))\) green
curves intersect \(\operatorname{cl}(\Delta')\).

\end{proof}

\end{document}
